%%%%%%%%%%%%%%%%%%%%%%%%%%%%%%%%%%%%%%%%%%%%%%%%%%%%%%%%%%%%%%%%%%%%%%%%
%% Art. 27, I — "precisar o setor técnico a que se aplica a invenção ou %%
%% o modelo de utilidade".                                              %%
%%                                                                     %%
%% Diga a qual setor técnico o seu objeto pertence: composições de      %%
%% tintura capilar, máquinas para semeadura, comunicações de rede sem   %%
%% fio, artigos domésticos... Se o objeto puder ser aplicado em mais de %%
%% um campo técnico, cite todos.                                        %%
%%                                                                     %%
%% Abra cada parágrafo com \pnum — ele produz a numeração sequencial    %%
%% [001], [002] exigida pelo art. 26, II.                               %%
%%%%%%%%%%%%%%%%%%%%%%%%%%%%%%%%%%%%%%%%%%%%%%%%%%%%%%%%%%%%%%%%%%%%%%%%

\secaoinpi{Campo \danatureza}

\pnum A presente \nomedanatureza\ pertence ao setor técnico dos equipamentos de
irrigação agrícola, mais especificamente ao dos dispositivos de acionamento de
válvulas empregadas no controle de vazão em sistemas de irrigação localizada.

\pnum A \nomedanatureza\ aplica-se também ao setor de automação de sistemas
hidráulicos de baixa pressão, no qual o acionamento de válvulas exige resposta
rápida e consumo reduzido de energia.
